% ======================================================================
\chapter*{Обозначения, единицы и физические постоянные}
\addcontentsline{toc}{chapter}{Обозначения, единицы и физические постоянные}
\markboth{Обозначения и постоянные}{}
% ======================================================================

\section*{Система единиц}\label{sec:units}

Во всём пособии используется \textbf{гауссова система единиц (СГС)}, традиционная для курса
общей физики МФТИ. В ней закон Кулона имеет вид $F=q_1q_2/r^2$, а магнитный момент тока
$\mu = IS/c$. Переход к СИ выполняется по правилам
\[
e^2 \;\longrightarrow\; \frac{e^2}{4\pi\varepsilon_0},\qquad
\frac{e}{c}\vect A \;\longrightarrow\; e\vect A,\qquad
\frac{e\hbar}{2\mel c}\;\longrightarrow\;\frac{e\hbar}{2\mel}.
\]
Например, $\aB=\hbar^2/(\mel e^2)$ в СГС соответствует $\aB = 4\pi\varepsilon_0\hbar^2/(\mel e^2)$ в СИ.
Численные значения приводятся в единицах, удобных для атомной и ядерной физики (эВ, Å, фм).
Заряд электрона обозначается $-e$, где $e>0$.

\section*{Основные обозначения}

\begin{center}
\small
\begin{tabularx}{\linewidth}{@{}>{\raggedright\arraybackslash}p{0.26\linewidth}X@{}}
\toprule
Обозначение & Смысл \\
\midrule
$\Psi(\vect r,t)$, $\psi(\vect r)$ & волновая функция; $\psi$ — её координатная часть для стационарного состояния \\
$\op A$ & оператор физической величины $A$ (шляпка отличает оператор от величины) \\
$\op{\vect p}=-\ii\hbar\nabla$ & оператор импульса \\
$\op H$ & гамильтониан (оператор полной энергии) \\
$\avg{A}=\int\Psi^*\op A\Psi\,\dd V$ & среднее значение величины $A$ в состоянии $\Psi$ \\
$\Delta A=\sqrt{\avg{A^2}-\avg{A}^2}$ & неопределённость (среднеквадратичное отклонение) \\
$\comm{\op A}{\op B}=\op A\op B-\op B\op A$ & коммутатор операторов \\
$\vect k$, $k=2\pi/\lambda$ & волновой вектор, волновое число \\
$\omega$, $\nu=\omega/2\pi$ & круговая и обычная частоты \\
$n,\ l,\ m$ ($m_l$), $m_s$ & главное, орбитальное, магнитное и спиновое квантовые числа \\
$\op{\vect L}$, $\op{\vect S}$, $\op{\vect J}$ & операторы орбитального, спинового и полного моментов (в единицах $\hbar$: $l,s,j$) \\
$\mu$ & приведённая масса (в атомной физике) или магнитный момент — смысл ясен из контекста \\
$Z,\ N,\ A$ & зарядовое число, число нейтронов, массовое число ядра \\
$\nucl{X}{A}$, $\nuc{X}{A}{Z}$ & нуклид с массовым числом $A$ и зарядом $Z$ \\
$\sh x,\ \ch x,\ \th x$ & гиперболические синус, косинус, тангенс \\
$D$, $R$ & коэффициенты прохождения и отражения \\
$\lambda$ & длина волны или постоянная распада — смысл ясен из контекста \\
\bottomrule
\end{tabularx}
\end{center}

\section*{Физические постоянные}

\begin{center}
\small
\begin{tabularx}{\linewidth}{@{}>{\raggedright\arraybackslash}p{0.38\linewidth}l>{\raggedright\arraybackslash}X@{}}
\toprule
Величина & Обозначение & Значение \\
\midrule
Скорость света & $c$ & $2{,}998\cdot10^{10}$ см/с \\
Постоянная Планка & $h$ & $6{,}626\cdot10^{-27}$ эрг$\cdot$с $=4{,}136\cdot10^{-15}$ эВ$\cdot$с \\
Приведённая постоянная Планка & $\hbar=h/2\pi$ & $1{,}0546\cdot10^{-27}$ эрг$\cdot$с $=6{,}582\cdot10^{-16}$ эВ$\cdot$с \\
Элементарный заряд & $e$ & $4{,}803\cdot10^{-10}$ ед. СГСЭ $=1{,}602\cdot10^{-19}$ Кл \\
Масса электрона & $\mel$ & $9{,}109\cdot10^{-28}$ г; $\mel c^2=0{,}511$ МэВ \\
Масса протона & $m_p$ & $1{,}6726\cdot10^{-24}$ г; $m_pc^2=938{,}27$ МэВ \\
Масса нейтрона & $m_n$ & $m_nc^2=939{,}57$ МэВ \\
Атомная единица массы & u & $1{,}6605\cdot10^{-24}$ г; $931{,}49$ МэВ$/c^2$ \\
Отношение масс & $m_p/\mel$ & $1836{,}15$ \\
Постоянная тонкой структуры & $\alpha=e^2/\hbar c$ & $1/137{,}036$ \\
Боровский радиус & $\aB=\hbar^2/\mel e^2$ & $0{,}5292\cdot10^{-8}$ см $=0{,}5292$ Å \\
Ридберг (энергия) & $\Ry=\mel e^4/2\hbar^2$ & $13{,}606$ эВ \\
Постоянная Ридберга & $R_\infty=\Ry/hc$ & $1{,}0974\cdot10^{5}$ см$^{-1}$ \\
Комптоновская длина волны электрона & $\lC=h/\mel c$ & $2{,}426\cdot10^{-10}$ см $=2{,}426$ пм \\
Магнетон Бора & $\muB=e\hbar/2\mel c$ & $9{,}274\cdot10^{-21}$ эрг/Гс $=5{,}788\cdot10^{-5}$ эВ/Тл \\
Ядерный магнетон & $\mu_N=e\hbar/2m_pc$ & $5{,}051\cdot10^{-24}$ эрг/Гс \\
Постоянная Больцмана & $k_{\mathrm B}$ & $1{,}381\cdot10^{-16}$ эрг/К $=8{,}617\cdot10^{-5}$ эВ/К \\
Комбинация & $\hbar c$ & $197{,}33$ МэВ$\cdot$фм $=1973{,}3$ эВ$\cdot$Å \\
Комбинация & $e^2$ & $1{,}440$ МэВ$\cdot$фм $=14{,}40$ эВ$\cdot$Å \\
Комбинация & $hc$ & $12\,398$ эВ$\cdot$Å \\
Постоянная Ферми & $G_F/(\hbar c)^3$ & $1{,}166\cdot10^{-5}$ ГэВ$^{-2}$ \\
\bottomrule
\end{tabularx}
\end{center}

\noindent Полезные соотношения: $1\eV = 1{,}602\cdot10^{-12}$ эрг; $1\fm=10^{-13}$ см;
$1\Ang=10^{-8}$ см; $1\,\text{Тл}=10^4$ Гс; энергия $1\eV$ соответствует температуре
$\approx 11\,600$ К и длине волны фотона $\lambda\approx 12\,398$~Å.
